%%% BEGIN fig:glauber-contour
\begin{figure}[htbp]
\centering
\begin{tikzpicture}[fd,
  t2 pole/.pic={\draw[line width=1pt,line cap=round] (-0.085,-0.085)--(0.085,0.085) (-0.085,0.085)--(0.085,-0.085);},
  t2 contour/.style={draw=blue!65!black,line width=1.15pt,postaction={decorate,decoration={markings,
      mark=at position #1 with {\arrow{Stealth[length=6pt,width=5pt]}}}}},
  t2 contour/.default=0.5,
  t2 ax/.style={lcaxis,gray!60!black},
  t2 band/.style={fill=red!9,draw=none},
  t2 lab/.style={font=\footnotesize}]
  % ======================= (a) all A-collinear poles on one side =======================
  \begin{scope}
    \fill[t2 band] (-0.55,-1.3) rectangle (0.55,1.45);
    \draw[red!45,line width=0.4pt,dash pattern=on 2pt off 1.5pt] (-0.55,-1.3) -- (-0.55,1.45) (0.55,-1.3) -- (0.55,1.45);
    \draw[t2 ax] (-2.4,0) -- (2.5,0) node[below,xshift=-6pt] {$\Real k^-$};
    \draw[t2 ax] (0,-1.3) -- (0,1.6) node[left] {$\Imag k^-$};
    \draw[t2 contour=0.5] (-2.3,0) -- (-1.75,0) .. controls (-1.4,0) and (-1.25,1.05) .. (0,1.05)
        .. controls (1.25,1.05) and (1.4,0) .. (1.75,0) -- (2.3,0);
    \pic at (-0.32,-0.17) {t2 pole}; \pic at (0.1,-0.27) {t2 pole}; \pic at (0.38,-0.12) {t2 pole};
    \node[t2 lab,anchor=east,align=right] at (-0.66,-0.62) {同侧的 $A$\\共线极点};
    \node[t2 lab,anchor=south west] at (0.9,0.95) {$|k^-|\sim|k_T|$};
    \node[t2 lab,text=red!55!black,anchor=south] at (0,1.62) {Glauber 区};
    \node at (0,-1.9) {(a)\ 极点同侧：围道可移出 Glauber 区};
  \end{scope}
  % ======================= (b) poles on both sides: Glauber pinch =======================
  \begin{scope}[xshift=6.4cm]
    \fill[t2 band] (-0.55,-1.3) rectangle (0.55,1.45);
    \draw[red!45,line width=0.4pt,dash pattern=on 2pt off 1.5pt] (-0.55,-1.3) -- (-0.55,1.45) (0.55,-1.3) -- (0.55,1.45);
    \draw[t2 ax] (-2.4,0) -- (2.5,0) node[below,xshift=-6pt] {$\Real k^-$};
    \draw[t2 ax] (0,-1.3) -- (0,1.6) node[left] {$\Imag k^-$};
    \draw[t2 contour=0.82] (-2.3,0) -- (2.3,0);
    \pic at (0.15,0.19) {t2 pole}; \pic at (-0.15,-0.19) {t2 pole};
    \draw[gray!60!black,line width=0.45pt] (0.24,0.27) -- (0.75,0.72) node[right,t2 lab,text=black,align=left] {活动部分子线\\（进入 $H$）};
    \draw[gray!60!black,line width=0.45pt] (-0.24,-0.27) -- (-0.75,-0.72) node[left,t2 lab,text=black] {旁观者线};
    \node[t2 lab,text=red!55!black,anchor=south] at (0,1.62) {Glauber 区};
    \node at (0,-1.9) {(b)\ 极点异侧：围道被钳在 Glauber 区};
  \end{scope}
\end{tikzpicture}
\caption{软胶子动量 $k$（固定 $k^+$、$k_T$，$k_T\sim\lambda Q$）在复 $k^-$ 平面上的 Glauber 区分析。阴影带为 Glauber 区 $|k^-|\lesssim k_T^2/Q$；
图中极点来自 $k$ 流经的 $A$ 共线传播子，位于 $|k^-|\sim k_T^2/Q$。约定 $k$ 在附着顶点从 $A$ 子图流出，并取末态旁观者动量 $p_X$ 固定：
Glauber 顶点之前的旁观者线 $l_s=p_X+k$ 中 $k$ 与大分量 $l_s^+$ 同向，极点在下半平面；进入 $H$ 的活动部分子线 $l_a=P_A-p_X-k$ 中 $k$ 与 $l_a^+$ 反向，极点在上半平面（$k$ 流入时上下互换）。
$B$ 共线线与胶子传播子本身给出的 $k^-$ 极点位于 $|k^-|\sim Q$，在图外。
(a)~所有 $A$ 共线极点在同一侧：围道可向另一侧变形到 $|k^-|\sim|k_T|$，在那里软近似成立；
(b)~旁观者线与活动部分子线的极点分居两侧：围道被钳在 Glauber 区内（Glauber 钳夹）。
仅有 $k^-$ 钳夹并不致命：若胶子另一端的线（如 DIS 中出射的被击出夸克）给出的 $k^+$ 极点都在同一侧，仍可把 $k^+$ 围道变形到 $|k^+|\sim|k_T|$，使该端的程函近似成立；
Drell--Yan 中 $B$ 一侧的 $k^+$ 围道同样被钳住，这一钳夹只有在对末态切割求和之后才被消除（\cref{fig:glauber-exchange}）。}
\label{fig:glauber-contour}
\end{figure}
%%% END fig:glauber-contour

%%% BEGIN fig:method-coeff-minus-coll
\begin{figure}[htbp]
\centering
% photons, photon vertices and the cut, common to all graphs (origin: bottom centre of a graph)
\newcommand{\fbCMCphot}{%
  \draw[photon] (-1.22,1.95) -- (-0.88,1.4);
  \draw[photon] (0.88,1.4) -- (1.22,1.95);
  \node[vertex] at (-0.88,1.4) {}; \node[vertex] at (0.88,1.4) {};
  \draw[cut] (0,-0.18) -- (0,2.02);}
% collinear approximation: operator vertices, Wilson line between them, Born struck quark on top;
% #1 = position of the arrow on the Wilson line
\newcommand{\fbCMCapprox}[1]{%
  \draw[quark] (-0.88,0.92) -- (-0.88,1.4);
  \draw[quark=0.78] (-0.88,1.4) -- (0.88,1.4);
  \draw[quark] (0.88,1.4) -- (0.88,0.92);
  \draw[wilson=#1] (-0.88,0.92) -- (0.88,0.92);
  \node[opvertex] at (-0.88,0.92) {}; \node[opvertex] at (0.88,0.92) {};
  \fbCMCphot}
\begin{tikzpicture}[fd,
  fb lab/.style={font=\footnotesize,align=center,anchor=north,inner sep=1pt},
  fb row/.style={font=\small,anchor=east},
  fb frame/.style={rounded corners=3pt,line width=0.6pt}]
  % ======================= row 1: one-loop DIS cut graphs (gamma* q) =======================
  % (i) ISR squared
  \begin{scope}[xshift=0cm]
    \draw[quark] (-0.88,0) -- (-0.88,0.58);
    \draw[quark] (-0.88,0.58) -- (-0.88,1.4);
    \draw[quark=0.78] (-0.88,1.4) -- (0.88,1.4);
    \draw[quark] (0.88,1.4) -- (0.88,0.58);
    \draw[quark] (0.88,0.58) -- (0.88,0);
    \draw[gluon] (-0.88,0.58) -- (0.88,0.58);
    \fbCMCphot
    \node[vertex] at (-0.88,0.58) {}; \node[vertex] at (0.88,0.58) {};
    \node[font=\scriptsize,below,inner sep=1.5pt] at (-0.88,0) {$p$};
    \node[font=\scriptsize,right,inner sep=1.5pt] at (-1.12,1.98) {$q$};
    \node[fb lab] at (0,-0.36) {初态辐射\\的平方};
  \end{scope}
  % (ii) ISR--FSR interference
  \begin{scope}[xshift=2.9cm]
    \draw[quark] (-0.88,0) -- (-0.88,0.58);
    \draw[quark] (-0.88,0.58) -- (-0.88,1.4);
    \draw[quark=0.3] (-0.88,1.4) -- (0.46,1.4);
    \draw[quark] (0.46,1.4) -- (0.88,1.4);
    \draw[quark] (0.88,1.4) -- (0.88,0);
    \draw[gluon] (-0.88,0.58) to[out=0,in=-90] (0.46,1.4);
    \fbCMCphot
    \node[vertex] at (-0.88,0.58) {}; \node[vertex] at (0.46,1.4) {};
    \node[fb lab] at (0,-0.36) {ISR--FSR 干涉\\（+镜像）};
  \end{scope}
  % (iii) FSR squared
  \begin{scope}[xshift=5.8cm]
    \draw[quark] (-0.88,0) -- (-0.88,1.4);
    \draw[quark] (-0.88,1.4) -- (-0.46,1.4);
    \draw[quark=0.27] (-0.46,1.4) -- (0.46,1.4);
    \draw[quark] (0.46,1.4) -- (0.88,1.4);
    \draw[quark] (0.88,1.4) -- (0.88,0);
    \draw[gluon] (-0.46,1.4) arc[start angle=180,end angle=0,radius=0.46];
    \fbCMCphot
    \node[vertex] at (-0.46,1.4) {}; \node[vertex] at (0.46,1.4) {};
    \node[fb lab] at (0,-0.36) {末态辐射\\的平方};
  \end{scope}
  % (iv) virtual vertex correction
  \begin{scope}[xshift=8.7cm]
    \draw[quark] (-0.88,0) -- (-0.88,0.58);
    \draw[quark] (-0.88,0.58) -- (-0.88,1.4);
    \draw[quark] (-0.88,1.4) -- (-0.35,1.4);
    \draw[quark=0.65] (-0.35,1.4) -- (0.88,1.4);
    \draw[quark] (0.88,1.4) -- (0.88,0);
    \draw[gluon] (-0.88,0.58) to[out=0,in=-90] (-0.35,1.4);
    \fbCMCphot
    \node[vertex] at (-0.88,0.58) {}; \node[vertex] at (-0.35,1.4) {};
    \node[fb lab] at (0,-0.36) {顶点虚修正\\（+镜像）};
  \end{scope}
  % ======================= row 2: collinear approximations T_A Gamma = (PDF graph) x Born =======================
  \begin{scope}[yshift=-4.05cm]
    % (i') ladder graph of the PDF operator
    \begin{scope}[xshift=0cm]
      \draw[quark] (-0.88,0) -- (-0.88,0.42);
      \draw[quark] (-0.88,0.42) -- (-0.88,0.92);
      \draw[quark] (0.88,0.92) -- (0.88,0.42);
      \draw[quark] (0.88,0.42) -- (0.88,0);
      \draw[gluon] (-0.88,0.42) -- (0.88,0.42);
      \fbCMCapprox{0.25}
      \node[vertex] at (-0.88,0.42) {}; \node[vertex] at (0.88,0.42) {};
      \node[font=\scriptsize,below,inner sep=1.5pt] at (-0.88,0) {$p$};
      \node[fb lab] at (0,-0.36) {PDF 算符的\\梯形图};
    \end{scope}
    % (ii') real Wilson-line graph
    \begin{scope}[xshift=2.9cm]
      \draw[quark] (-0.88,0) -- (-0.88,0.42);
      \draw[quark] (-0.88,0.42) -- (-0.88,0.92);
      \draw[quark] (0.88,0.92) -- (0.88,0);
      \draw[gluon] (-0.88,0.42) to[out=0,in=-90] (0.46,0.92);
      \fbCMCapprox{0.22}
      \node[vertex] at (-0.88,0.42) {}; \node[vertex] at (0.46,0.92) {};
      \node[fb lab] at (0,-0.36) {Wilson 线实图\\（+镜像）};
    \end{scope}
    % (iii') no target-collinear region
    \begin{scope}[xshift=5.8cm]
      \node[font=\large] at (0,0.9) {$0$};
      \node[fb lab] at (0,-0.36) {无靶共线区};
    \end{scope}
    % (iv') virtual Wilson-line graph
    \begin{scope}[xshift=8.7cm]
      \draw[quark] (-0.88,0) -- (-0.88,0.42);
      \draw[quark] (-0.88,0.42) -- (-0.88,0.92);
      \draw[quark] (0.88,0.92) -- (0.88,0);
      \draw[gluon] (-0.88,0.42) to[out=0,in=-90] (-0.35,0.92);
      \fbCMCapprox{0.72}
      \node[vertex] at (-0.88,0.42) {}; \node[vertex] at (-0.35,0.92) {};
      \node[fb lab] at (0,-0.36) {Wilson 线虚图\\（+镜像）};
    \end{scope}
  \end{scope}
  % ======================= colour frames, row labels, result =======================
  \begin{scope}[on background layer]
    \foreach \x/\c in {0/green,2.9/blue,5.8/gray,8.7/orange}
      \draw[fb frame,draw=\c!55!black,fill=\c!6] (\x-1.4,-5.3) rectangle (\x+1.4,2.2);
  \end{scope}
  \node[fb row] at (-1.5,0.9) {$\Gamma^{(1)}$};
  \node[fb row] at (-1.5,-3.15) {$-\,T_A\Gamma^{(1)}$};
  \foreach \x in {0,2.9,5.8,8.7} \node[font=\large] at (\x,-1.62) {$-$};
  \draw[decorate,decoration={brace,amplitude=5pt},line width=0.6pt] (10.25,2.15) -- (10.25,-5.25);
  \node[anchor=west,align=left,font=\small] at (10.45,-1.55) {各列之和\\[3pt]$=C_H\Gamma^{(1)}$\\[3pt]$=\dfrac{\as}{2\pi}\,C_2^q$\\[1pt]{\footnotesize（$\mu=Q$）}};
\end{tikzpicture}
\caption{单圈 DIS 系数函数 = 割图减去其共线近似（$\gamma^*q$ 道，Feynman 规范；竖直红色虚线为末态割线，左侧为振幅）。
上行为 $\order{\as}$ 割图 $\Gamma^{(1)}$，下行为其靶共线近似 $T_A\Gamma^{(1)}=f^{(1)}_{q/q}\otimes{\rm Born}$：PDF 算符的割图（$\otimes$ 为算符顶点，双线为沿 $n_B$ 的 Wilson 线）乘以领头阶硬部分；同色框中的图一一对应。
初态辐射的平方对应 PDF 算符的梯形图，ISR--FSR 干涉对应胶子连到割线另一侧 Wilson 线的实图，顶点虚修正对应同侧的 Wilson 线虚图（外线自能无标度，未画出）；末态辐射没有靶共线区，不需减除。
各列之差的和即硬区减除项 $C_H\Gamma^{(1)}$，在 $\mu=Q$ 时为 $\frac{\as}{2\pi}C_2^q$；纯维数正规化中下行各图无标度，只留下 $\MSbar$ 抵消项，$C_A\Gamma^{(1)}=-\frac{\as}{2\pi}\frac1\varepsilon P_{qq}$（VF-DIS-6，VF-PDF-3）。}
\label{fig:method-coeff-minus-coll}
\end{figure}
%%% END fig:method-coeff-minus-coll

%%% BEGIN fig:method-flowchart
\begin{figure}[htbp]
\centering
\newcommand{\fbref}[1]{\mbox{\cref{#1}}}
\begin{tikzpicture}[fd,font=\footnotesize,
  fb num/.style={circle,fill=gray!65!black,text=white,font=\scriptsize\bfseries,inner sep=0pt,minimum size=0.42cm},
  fb title/.style={font=\small\bfseries,anchor=west,inner sep=1pt},
  fb stitle/.style={font=\footnotesize\bfseries,anchor=west,inner sep=1pt},
  fb cell/.style={rectangle,rounded corners=1.5pt,line width=0.5pt,text width=3.38cm,minimum height=1.0cm,
                  align=left,inner sep=3pt,font=\footnotesize},
  fb R/.style={fb cell,fill=OliveGreen!10,draw=OliveGreen!65!black},
  fb Rc/.style={fb R,dash pattern=on 2.4pt off 1.4pt},
  fb A/.style={fb cell,fill=RoyalBlue!8,draw=RoyalBlue!65!black},
  fb frame/.style={draw=gray!55!black,line width=0.65pt,rounded corners=3pt,fill=white,inner sep=3pt},
  fb sub/.style={draw=gray!45,line width=0.45pt,rounded corners=2pt,fill=gray!4,inner sep=2pt},
  fb htext/.style={text width=6.45cm,align=left,inner sep=1pt,font=\footnotesize},
  fb H/.style={rectangle,rounded corners=2pt,draw=orange!70!black,line width=0.55pt,dash pattern=on 2.4pt off 1.4pt,
               fill=orange!5,inner sep=3pt},
  fb link/.style={draw=orange!70!black,line width=0.5pt,dash pattern=on 2.4pt off 1.4pt},
  fb flow/.style={draw=gray!55!black,line width=1.1pt,-{Stealth[length=6pt,width=5pt]}}]
  % one step: #1 name, #2 y of its title, #3 number, #4 title, #5 R cell, #6 A cell, #7 R style
  \newcommand{\fbstep}[7]{%
    \node[fb num] (#1n) at (0.28,#2) {#3};
    \node[fb title] (#1t) at ($(#1n.east)+(0.06,0)$) {#4};
    \node[#7,anchor=north west] (#1r) at (0.12,#2-0.3) {#5};
    \node[fb A,anchor=north west] (#1a) at ($(#1r.north east)+(0.12,0)$) {#6};
    \begin{scope}[on background layer]\node[fb frame,fit=(#1n)(#1t)(#1r)(#1a)] (#1) {};\end{scope}}
  % ---------------- steps 1--4 ----------------
  \fbstep{s1}{0}{1}{区域：Landau 方程与幂次计数}
    {幂次计数（标度陈述）\\\fbref{prop:ir-power-counting}}
    {领头区域与约化图\\\fbref{thm:ir-leading-regions}}{fb R}
  \fbstep{s2}{-2.1}{2}{近似：共线与软近似子（沿用原有 $\ii0$）}
    {$\ii0$ 相容性\\\fbref{lem:method-i0-compat}}
    {近似精度\\\fbref{lem:method-collinear-accuracy,lem:method-soft-accuracy}}{fb R}
  \fbstep{s3}{-4.2}{3}{Grammer--Yennie 分解\enspace$g_{\mu\nu}=K_{\mu\nu}+G_{\mu\nu}$}
    {分解恒等式\\\fbref{lem:method-gy}}
    {G 部分幂次压低\\\fbref{prop:method-gy-suppression}}{fb R}
  \fbstep{s4}{-6.3}{4}{Ward 恒等式 $\to$ Wilson 线}
    {单个与树图两个 K 胶子\\\fbref{prop:method-ward-one,prop:method-ward-tree}}
    {多胶子（含鬼场）\\\fbref{lem:method-ward-multi}}{fb R}
  % ---------------- step 5: Glauber region, unitarity, subtraction ----------------
  \node[fb num] (s5n) at (0.28,-8.4) {5};
  \node[fb title] (s5t) at ($(s5n.east)+(0.06,0)$) {Glauber 区、幺正性与减除};
  \foreach \k/\y/\ttl/\rc/\ac/\rs in {%
      a/-8.85/{5a\enspace Glauber 区：沿 $e_G=n_B-n_A$ 同时形变}/{单胶子交换的极点结构\\\fbref{prop:method-glauber-oneloop}}/{围道形变判据\\\fbref{prop:method-glauber-deformation}}/{fb R},
      b/-10.75/{5b\enspace 幺正性：末态相互作用与软因子}/{软因子的幺正性\\\fbref{lem:method-soft-unitarity}}/{末态相互作用相消\\\fbref{prop:method-fsi}}/{fb R},
      c/-12.65/{5c\enspace 减除（Collins）：$\Gamma=\sum_RC_R\Gamma+r$}/{链展开与乘积公式\\\fbref{prop:method-chains,prop:method-product}}/{减除定理\\\fbref{thm:method-subtraction}}/{fb R}} {
    \node[fb stitle] (g5\k t) at (0.22,\y) {\ttl};
    \node[\rs,anchor=north west,text width=3.3cm] (g5\k r) at (0.22,\y-0.24) {\rc};
    \node[fb A,anchor=north west,text width=3.3cm] (g5\k a) at ($(g5\k r.north east)+(0.12,0)$) {\ac};
    \begin{scope}[on background layer]\node[fb sub,fit=(g5\k t)(g5\k r)(g5\k a)] (s5\k) {};\end{scope}
  }
  \begin{scope}[on background layer]
    \node[fb frame,fit=(s5n)(s5t)(s5a)(s5b)(s5c)] (s5) {};
  \end{scope}
  % ---------------- step 6 ----------------
  \node[fb num] (s6n) at (0.28,-15.05) {6};
  \node[fb title] (s6t) at ($(s6n.east)+(0.06,0)$) {因子化形式};
  \node[fb A,anchor=north west,text width=7.09cm,minimum height=0.62cm,align=center] (s6a) at (0.12,-15.35)
    {$\sigma=\sum_{ij}H_{ij}\otimes f_{i/A}\otimes f_{j/B}\,[\otimes S]+\order{\lambda^p}$};
  \begin{scope}[on background layer]\node[fb frame,fit=(s6n)(s6t)(s6a)] (s6) {};\end{scope}
  % ---------------- flow arrows ----------------
  \foreach \u/\v in {s1/s2,s2/s3,s3/s4,s4/s5,s5/s6} \draw[fb flow] (\u.south) -- (\v.north);
  % ---------------- hypotheses (H1)--(H5) with the numbered assumptions ----------------
  \node[fb htext,anchor=west] (h1t) at (8.2,0|-s1) {\textbf{(H1)} 领头区域；积分范数下的幂次计数\\\fbref{asm:appN-envelope,asm:appN-separation}};
  \node[fb htext,anchor=west] (h2t) at (8.2,0|-s3) {\textbf{(H2)} $H$--$J$ 界面上共线近似成立，K 胶子退耦为 Wilson 线\\\fbref{asm:appN-ward-product}\\[1pt]
      {\color{gray!35!black}（步骤 2--4 中的软部分属于 (H3)）}};
  \node[fb htext,anchor=north west] (h3t) at ($(8.2,0|-s5a)+(0,0.66)$) {\textbf{(H3)} Glauber 区之后：软近似与软胶子退耦\\\fbref{asm:appO-mixing,asm:appO-deformation,asm:appO-rewritten-cs}};
  \node[fb Rc,anchor=north west,text width=6.3cm,minimum height=0.48cm] (h3c) at ($(h3t.south west)+(0,-0.07)$) {DY 的精确形式：\fbref{thm:appO-main}（$R_{\rm cond}$）};
  \node[fb htext,anchor=west] (h4t) at (8.2,0|-s5b) {\textbf{(H4)} 末态相互作用由幺正性相消；软因子为 $1$ 或 Wilson 线真空矩阵元\\\fbref{asm:appN-cutsum,asm:appO-lcopt}};
  \node[fb htext,anchor=north west] (h5t) at ($(8.2,0|-s5c)+(0,0.66)$) {\textbf{(H5)} 减除项 = 重整化的算符矩阵元\\\fbref{asm:appN-collins,asm:appN-values,asm:appN-uv,asm:appN-multren}};
  \node[fb Rc,anchor=north west,text width=6.3cm,minimum height=0.48cm] (h5c) at ($(h5t.south west)+(0,-0.07)$) {精确演绎：\fbref{thm:appN-subtraction}（$R_{\rm cond}$）};
  \begin{scope}[on background layer]
    \node[fb H,fit=(h1t)] (h1) {};
    \node[fb H,fit=(h2t)] (h2) {};
    \node[fb H,fit=(h3t)(h3c)] (h3) {};
    \node[fb H,fit=(h4t)] (h4) {};
    \node[fb H,fit=(h5t)(h5c)] (h5) {};
  \end{scope}
  \draw[fb link] (s1.east) -- (h1.west|-s1.east);
  \draw[fb link] (s2.east) -- ++(0.25,0) |- (h2.west);
  \draw[fb link] (s3.east) -- (h2.west|-s3.east);
  \draw[fb link] (s4.east) -- ++(0.25,0) |- (h2.west);
  \draw[fb link] (s5a.east) -- (h3.west|-s5a.east);
  \draw[fb link] (s5b.east) -- (h4.west|-s5b.east);
  \draw[fb link] (s5c.east) -- (h5.west|-s5c.east);
  \node[font=\footnotesize,anchor=west,align=left] at (8.2,0|-s6) {$\Longleftarrow$\ \mbox{(H1)--(H5)} 成立时的结论：\\\fbref{thm:method-template}（模板定理），全阶论证 [A]};
  % ---------------- legend ----------------
  \node[fb R,text width=0.3cm,minimum height=0.3cm,inner sep=0pt,anchor=west] (l1) at (0.1,-16.55) {};
  \node[anchor=west] (l1t) at (l1.east) {[R]（在所用层次上严格；含数学引理）};
  \node[fb A,text width=0.3cm,minimum height=0.3cm,inner sep=0pt,anchor=west] (l2) at ($(l1t.east)+(0.35,0)$) {};
  \node[anchor=west] (l2t) at (l2.east) {[A] 全阶论证};
  \node[fb Rc,text width=0.3cm,minimum height=0.3cm,inner sep=0pt,anchor=west] (l3) at ($(l2t.east)+(0.35,0)$) {};
  \node[anchor=west] (l3t) at (l3.east) {条件定理 $R_{\rm cond}$};
  \node[fb H,minimum width=0.36cm,minimum height=0.32cm,inner sep=0pt,anchor=west] (l4) at ($(l3t.east)+(0.35,0)$) {};
  \node[anchor=west] (l4t) at (l4.east) {编号假设};
\end{tikzpicture}
\caption{全阶因子化证明的六个步骤（左），它们与\cref{thm:method-template}（模板定理）的假设 \mbox{(H1)--(H5)} 的对应，以及精确形式中的编号假设（右；\cref{app:appN}与\cref{app:appO}，对应关系即\cref{tab:method-hypotheses}）。
每个步骤中，绿色框为在所用层次上已严格证明的陈述（[R] 或数学引理），蓝色框为全阶论证 [A]，绿色虚线框为条件定理 $R_{\rm cond}$，即\cref{app:appN}中的减除演绎（\cref{thm:appN-subtraction}）与\cref{app:appO}中单举色单态 Drell--Yan 过程的 Glauber 相消（\cref{thm:appO-main}）。
步骤 4 中，$A$ 共线的 K 胶子化为沿 $n_B$ 的 Wilson 线（$B$ 共线的则沿 $n_A$），软胶子化为沿各喷注方向的 Wilson 线；离开 Glauber 区须沿 $e_G=n_B-n_A$ 同时形变两个光锥分量（\cref{fig:method-glauber-direction}），软因子的幺正性只在处理完 Glauber 区之后使用。}
\label{fig:method-flowchart}
\end{figure}
%%% END fig:method-flowchart

%%% BEGIN fig:method-glauber-direction
\begin{figure}[htbp]
\centering
% real (k^+,k^-) plane with the Glauber wedge G_delta = {|k^+ k^-| <= delta k_T^2}; units |k_T| = 1, delta = 0.12
\newcommand{\fbGwedge}{%
  \fill[red!9] plot[domain=1.6:0.075,samples=40] (\x,{0.12/\x})
      -- plot[domain=0.075:1.6,samples=40] ({-\x},{0.12/\x})
      -- plot[domain=1.6:0.075,samples=40] ({-\x},{-0.12/\x})
      -- plot[domain=0.075:1.6,samples=40] (\x,{-0.12/\x}) -- cycle;
  \foreach \s/\t in {1/1,-1/1,1/-1,-1/-1}
    \draw[red!55!black,line width=0.5pt] plot[domain=0.075:1.6,samples=40] ({\s*\x},{\t*0.12/\x});
  \draw[lcaxis] (-1.7,0) -- (1.8,0) node[right=0pt,font=\footnotesize] {$k^+$};
  \draw[lcaxis] (0,-1.7) -- (0,1.8) node[above=0pt,font=\footnotesize] {$k^-$};
  \draw[red!50!black,line width=0.4pt] (1.0,-0.06) -- (1.22,-0.45);
  \node[font=\footnotesize,text=red!50!black,anchor=north,inner sep=1pt] at (1.25,-0.45) {$G_\delta$};
  \fill (0.07,0.07) circle[radius=1.6pt];}
\begin{tikzpicture}[fd,
  fb lab/.style={font=\footnotesize,inner sep=1.5pt},
  fb text/.style={font=\footnotesize,align=center,anchor=north,inner sep=1pt},
  fb pole/.pic={\draw[line width=0.9pt,line cap=round] (-0.075,-0.075)--(0.075,0.075) (-0.075,0.075)--(0.075,-0.075);},
  fb def/.style={line width=1.3pt,-{Stealth[length=6pt,width=5pt]}},
  fb contour/.style={draw=blue!60!black,line width=1.0pt},
  fb carrow/.style={postaction={decorate,decoration={markings,
      mark=at position #1 with {\arrow{Stealth[length=5.5pt,width=4.5pt]}}}}}]
  % ======================= (a) one-component deformation stays in G_delta =======================
  \begin{scope}[x=1.15cm,y=1.15cm]
    \fbGwedge
    \draw[fb def,red!65!black] (0.07,0.12) -- (0.07,1.3);
    \node[fb lab,anchor=west,text=red!55!black] at (0.14,1.3) {$\Imag k^-\sim|\mathbf k_T|$};
    \node[font=\small,anchor=south west] at (-1.75,1.6) {(a)};
    \node[fb text] at (0,-1.87) {只形变 $k^-$（$\kappa\lesssim|\mathbf k_T|$）：\\$|k^+k^-|\sim\lambda^3\ll k_T^2$，\\仍在 $G_\delta$ 中};
  \end{scope}
  % ======================= (b) the soft line's own pole limits a one-component deformation =======================
  \begin{scope}[xshift=5.15cm]
    \draw[lcaxis] (-2.1,0) -- (2.1,0) node[right=0pt,font=\footnotesize] {$\Real k^-$};
    \draw[lcaxis] (0,-1.2) -- (0,1.75) node[above=0pt,font=\footnotesize] {$\Imag k^-$};
    \draw[fb contour,fb carrow=0.08,fb carrow=0.92] (-2.05,0) -- (-1.0,0) to[out=0,in=180] (0.55,0.78)
        to[out=0,in=180] (1.75,0) -- (2.05,0);
    \pic at (0.45,-0.16) {fb pole}; \pic at (0.8,-0.16) {fb pole}; \pic at (1.15,-0.16) {fb pole};
    \node[fb lab,anchor=north] at (0.85,-0.25) {$J_A$ 的极点};
    \pic[red!60!black] at (-1.45,0.16) {fb pole};
    \node[fb lab,anchor=south,text=red!55!black] at (-1.45,0.26) {软线极点};
    \draw[{Stealth[length=3.6pt,width=3pt]}-{Stealth[length=3.6pt,width=3pt]},line width=0.5pt] (0.55,0.04) -- (0.55,0.74);
    \node[fb lab,anchor=west] at (0.59,0.4) {$\kappa$};
    \node[font=\small,anchor=south west] at (-2.1,1.38) {(b)};
    \node[fb text] at (0,-2.15) {$k^+<0$：软线极点 $k^-=\frac{k_T^2-\ii0}{2k^+}$\\在上半平面，只形变 $k^-$ 的\\幅度 $\kappa\lesssim k_T^2/|k^+|$};
  \end{scope}
  % ======================= (c) simultaneous deformation along e_G leaves G_delta =======================
  \begin{scope}[xshift=10.35cm,x=1.15cm,y=1.15cm]
    \fbGwedge
    \draw[fb def,green!45!black] (0.03,0.11) -- (-0.92,1.06);
    \node[fb lab,anchor=north east,text=green!30!black] at (-0.5,0.5) {$e_G$};
    \node[font=\small,anchor=south west] at (-1.75,1.6) {(c)};
    \node[fb text] at (0,-1.87) {沿 $e_G=n_B-n_A$ 同时形变：\\$\Imag k^->0$，$\Imag k^+<0$，\\幅度 $\sim|\mathbf k_T|$：离开 $G_\delta$};
  \end{scope}
  % identity below the panels
  \node[font=\footnotesize,anchor=north] at (5.15,-3.82) {所有分母更类时：\ $(l+\ii t\sigma f e_G)^2-m^2=(l^2-m^2+2t^2f^2)+2\ii t\sigma f\,(l^+-l^-)$，\quad $e_G=(-1^+,+1^-,\mathbf 0_T)$};
\end{tikzpicture}
\caption{离开 Glauber 区需要沿 $e_G=n_B-n_A$ 同时形变两个光锥分量（示意；(a)、(c) 以 $|\mathbf k_T|$ 为单位，红色阴影为 Glauber 楔 $G_\delta=\{|k^+k^-|\le\delta k_T^2\}$，复形变量画在同一平面中）。
(a)~对 Glauber 动量 $k^\pm\sim\lambda^2$、$k_T\sim\lambda$ 只形变 $k^-$，即使形变到 $|\mathbf k_T|$ 量级，仍有 $|k^+k^-|\sim\lambda^3\ll k_T^2$，动量留在 $G_\delta$ 中；
(b)~$k^+<0$ 时软线自身的极点位于上半平面，与 $J_A$ 线的极点分居实轴两侧，把只形变 $k^-$ 的幅度限制在 $\kappa\sim k_T^2/|k^+|$（\cref{lem:method-i0-compat}）；
(c)~振幅一侧（$\sigma=+1$）流入 $J_A$ 的软动量沿 $e_G$ 形变（$\Imag k^->0$、$\Imag k^+<0$），幅度 $\sim|\mathbf k_T|$，离开 $G_\delta$，且使所有分母更类时（VF-REF-7；\cref{prop:method-glauber-deformation}）。}
\label{fig:method-glauber-direction}
\end{figure}
%%% END fig:method-glauber-direction

%%% BEGIN fig:method-hasse
\begin{figure}[htbp]
\centering
\begin{tikzpicture}[fd,
  fb J1/.style={draw=blue!60!black},
  fb J2/.style={draw=orange!80!black},
  fb S/.style={draw=green!45!black},
  fb Hl/.style={draw=black,line width=1.6pt},
  fb m/.style={font=\scriptsize,inner sep=2.6pt,sloped},
  fb node/.style={draw=gray!60!black,fill=white,rounded corners=2.5pt,line width=0.6pt,inner sep=2.5pt,
                  minimum width=1.5cm,minimum height=0.55cm,font=\small},
  fb edge/.style={draw=gray!50!black,line width=0.7pt},
  fb hl/.style={draw=red!60!black,line width=1.3pt}]
  \newcommand{\fbH}{H}
  \newcommand{\fbA}{{\color{blue!60!black}C_1}}
  \newcommand{\fbB}{{\color{orange!80!black}C_2}}
  \newcommand{\fbS}{{\color{green!45!black}S}}
  % ======================= left: two-loop planar vertex, coloured for the region (C1,S) =======================
  \begin{scope}
    \coordinate (V) at (0,0);
    \coordinate (Q1) at (28:1.9); \coordinate (Q2) at (28:3.4); \coordinate (E1) at (28:4.55);
    \coordinate (A1) at (-28:1.9); \coordinate (A2) at (-28:3.4); \coordinate (E2) at (-28:4.55);
    \draw[photon] (-1.3,0) -- (V);
    \draw[quark=0.82,fb Hl] (A1) -- node[fb m,below,pos=0.4,text=black] {$p_2+k_1+k_2$} (V);
    \draw[quark=0.2,fb J1] (V) -- node[fb m,above,pos=0.6,text=blue!50!black] {$p_1-k_1-k_2$} (Q1);
    \draw[quark=0.22,fb J1] (Q1) -- node[fb m,above,pos=0.62,text=blue!50!black] {$p_1-k_2$} (Q2);
    \draw[quark] (Q2) -- (E1);
    \draw[quark=0.8,fb J2] (A2) -- node[fb m,below,pos=0.38,text=orange!60!black] {$p_2+k_2$} (A1);
    \draw[quark] (E2) -- (A2);
    \draw[gluon,fb J1] (A1) -- (Q1);
    \draw[gluon,fb S] (A2) -- (Q2);
    \foreach \p in {V,Q1,Q2,A1,A2} \node[vertex] at (\p) {};
    \path[mom={$q$},mom sep=0.28cm] (-1.3,0) -- (V);
    \path[mom={$p_1$},mom pos=0.6,mom length=0.45cm] (Q2) -- (E1);
    \path[mom={$p_2$},mom pos=0.6,mom length=0.45cm] (A2) -- (E2);
    \path[mom'={\scriptsize$k_1$},mom length=0.42cm,mom sep=0.3cm] (A1) -- (Q1);
    \path[mom'={\scriptsize$k_2$},mom length=0.42cm,mom sep=0.3cm] (A2) -- (Q2);
    \node[font=\footnotesize,align=center] at (1.75,-2.75) {例：区域 $(R(k_1),R(k_2))=(\fbA,\fbS)$\\[1pt]
      {\color{blue!60!black}$J_1$}\quad{\color{orange!80!black}$J_2$}\quad{\color{green!45!black}$S$}\quad 粗线：$H$};
  \end{scope}
  % ======================= right: Hasse diagram of the nine leading regions =======================
  \begin{scope}[xshift=8.6cm,yshift=-2.35cm]
    \node[fb node] (SS) at (0,0) {$(\fbS,\fbS)$};
    \node[fb node] (AS) at (-1.45,1.15) {$(\fbA,\fbS)$};
    \node[fb node] (BS) at (1.45,1.15) {$(\fbB,\fbS)$};
    \node[fb node] (AA) at (-2.75,2.3) {$(\fbA,\fbA)$};
    \node[fb node] (HS) at (0,2.3) {$(\fbH,\fbS)$};
    \node[fb node] (BB) at (2.75,2.3) {$(\fbB,\fbB)$};
    \node[fb node] (HA) at (-1.45,3.45) {$(\fbH,\fbA)$};
    \node[fb node] (HB) at (1.45,3.45) {$(\fbH,\fbB)$};
    \node[fb node] (HH) at (0,4.6) {$(\fbH,\fbH)$};
    \foreach \u/\v in {SS/AS,SS/BS,AS/AA,AS/HS,BS/BB,BS/HS,HS/HA,HS/HB,BB/HB,HA/HH,HB/HH}
      \draw[fb edge] (\u) -- (\v);
    \draw[fb hl] (AA) -- (HA);
    \node[font=\scriptsize,text=red!55!black,align=right,anchor=east] at (-2.35,3.5) {逐线比较：\\$C_1\le H$，$C_1\le C_1$};
    % beta_R column
    \node[font=\footnotesize,text=gray!30!black] at (4.25,5.35) {$\beta_R$};
    \foreach \y/\b in {0/{1/81},1.15/{1/27},2.3/{1/9},3.45/{1/3}}
      \node[font=\footnotesize,text=gray!30!black] at (4.25,\y) {$\b$};
    \node[font=\scriptsize,text=gray!30!black] at (4.25,4.6) {$N=\Omega$};
    % pinch sets grow upwards
    \draw[lcaxis,gray!50!black,line width=0.7pt] (5.05,-0.2) -- (5.05,4.8);
    \node[font=\footnotesize,text=gray!30!black,rotate=90,anchor=north] at (5.12,2.3) {钳夹集增大};
    \node[font=\footnotesize,anchor=south] at (0,5.1) {$(R(k_1),R(k_2))$：内胶子与外胶子的模式};
  \end{scope}
\end{tikzpicture}
\caption{两圈平面顶点九个领头区域的 Hasse 图。左：内胶子 $k_1$、外胶子 $k_2$ 的动量走向（胶子动量由反夸克线流向夸克线），线的颜色给出区域 $(C_1,S)$ 的逐线模式：
$J_1=\{k_1,\,p_1-k_2,\,p_1-k_1-k_2\}$、$J_2=\{p_2+k_2\}$、$S=\{k_2\}$，硬线 $p_2+k_1+k_2$ 与光子顶点组成 $H$。
右：16 种模式组合中恰有 9 个领头钳夹，按逐线偏序（$S<C_i<H$，$C_1\wedge C_2=S$）排成 Hasse 图（12 条覆盖关系，自下而上钳夹集增大），其中 $(C_1,C_1)<(H,C_1)$（红线）不易一眼看出；
这一集合对逐线交运算封闭，例如 $(C_1,C_1)\wedge(H,S)=(C_1,S)$、$(H,C_1)\wedge(H,C_2)=(H,S)$、$(C_1,C_1)\wedge(C_2,C_2)=(S,S)$。
右列为各层的邻域指数 $\beta_R=3^{h(R)-h_*}$（$h_*=4$；最大区域的邻域为整个积分域 $\Omega$，\cref{def:method-neighbourhoods}）。}
\label{fig:method-hasse}
\end{figure}
%%% END fig:method-hasse

%%% BEGIN fig:method-kg
\begin{figure}[htbp]
\centering
% collinear blob A, hard blob, photon and outgoing struck quark (origin at the centre of A); #1 = name of the hard blob
\newcommand{\fbKGframe}[1]{%
  \node[jet blob,ellipse,minimum width=1.75cm,minimum height=0.66cm,inner sep=0pt] (A) at (0,0) {};
  \node[font=\small] at (-0.42,0) {$A$};
  \node[hard blob,ellipse,minimum width=1.6cm,minimum height=0.66cm,inner sep=0pt] (H) at (0,2.25) {};
  \node[font=\small] at (H) {$#1$};
  \draw[photon] (-1.25,3.05) -- (H.150);
  \draw[quark] (H.20) -- (1.25,3.0);}
\begin{tikzpicture}[fd,
  fb lab/.style={font=\footnotesize},
  fb idx/.style={font=\scriptsize,inner sep=1pt},
  fb op/.style={font=\large},
  fb box/.style={draw=gray!60!black,line width=0.5pt,rounded corners=3pt,fill=gray!4,inner sep=4pt}]
  % ======================= row 1: g = K + G =======================
  % full propagator numerator
  \begin{scope}
    \fbKGframe{H}
    \draw[quark] (-0.45,0.3) -- (-0.45,1.96);
    \draw[gluon] (0.42,0.33) -- (0.42,1.94);
    \path[mom'={$k$},mom sep=0.32cm] (0.42,0.33) -- (0.42,1.94);
    \node[fb idx,anchor=east] at (0.33,1.9) {$\mu$};
    \node[fb idx,anchor=east] at (0.34,0.5) {$\nu$};
    \node[fb lab] at (0,-0.75) {$g_{\mu\nu}$};
  \end{scope}
  \node[fb op] at (1.85,1.15) {$=$};
  % K part: scalar polarization k_mu at the hard end (dot)
  \begin{scope}[xshift=3.7cm]
    \fbKGframe{H}
    \draw[quark] (-0.45,0.3) -- (-0.45,1.96);
    \draw[longit] (0.42,0.33) -- (0.42,1.94);
    \fill[blue!60!black] (0.42,1.76) circle[radius=1.7pt];
    \path[mom'={$k$},mom sep=0.32cm] (0.42,0.33) -- (0.42,1.94);
    \node[fb idx,anchor=east] at (0.33,1.9) {$\mu$};
    \node[fb idx,anchor=east] at (0.34,0.5) {$\nu$};
    \node[fb lab] at (0,-0.85) {$K_{\mu\nu}=\dfrac{k_\mu\,n_{B\nu}}{n_B\cdot k-\ii0}$};
  \end{scope}
  \node[fb op] at (5.62,1.15) {$+$};
  % G part
  \begin{scope}[xshift=7.55cm]
    \fbKGframe{H}
    \draw[quark] (-0.45,0.3) -- (-0.45,1.96);
    \draw[gluon] (0.42,0.33) -- (0.42,1.94);
    \path[mom'={$k$},mom sep=0.32cm] (0.42,0.33) -- (0.42,1.94);
    \node[fb idx,anchor=east] at (0.33,1.9) {$\mu$};
    \node[fb idx,anchor=east] at (0.34,0.5) {$\nu$};
    \node[fb lab] at (0,-0.75) {$G_{\mu\nu}=g_{\mu\nu}-K_{\mu\nu}$};
    \node[fb lab] at (0,-1.28) {$H\cdot G\cdot C=H\cdot\big(C-\frac{C^+}{k^+}\,k\big)$};
    \node[fb lab,text=red!55!black] at (0,-1.78) {$\ge\lambda$ 压低};
  \end{scope}
  % ======================= row 2: Ward identity for the K part; the exception =======================
  \draw[draw=gray!50!black,line width=1.0pt,-{Stealth[length=6pt,width=5pt]}] (3.35,-1.42) -- (2.3,-2.5);
  \node[fb lab,anchor=east,align=right] at (2.42,-1.86) {对 $H$ 内一切附着求和\\（Ward 恒等式，\cref{fig:ward-identity}）};
  \begin{scope}[xshift=1.6cm,yshift=-5.85cm]
    \node[jet blob,ellipse,minimum width=1.75cm,minimum height=0.66cm,inner sep=0pt] (A) at (0,0) {};
    \node[font=\small] at (-0.42,0) {$A$};
    \node[hard blob,ellipse,minimum width=1.6cm,minimum height=0.66cm,inner sep=0pt] (H) at (0,2.25) {};
    \node[font=\small] at (H) {$H_0$};
    \draw[photon] (-1.25,3.05) -- (H.150);
    \draw[quark] (H.20) -- (1.25,3.0);
    \draw[quark] (-0.45,0.3) -- (H.236);
    \draw[wilson=0.8] (H.236) -- ++(-24:2.45) coordinate (W);
    \node[fb lab,anchor=west] at (W) {沿 $n_B$ 至 $\infty$};
    \coordinate (Y) at ($(H.236)!0.47!(W)$);
    \draw[longit] (0.42,0.33) -- (Y);
    \node[vertex] at (Y) {};
    \path[mom'={$k$},mom sep=0.3cm,mom pos=0.45] (0.42,0.33) -- (Y);
    \node[fb lab,anchor=north west,inner sep=1pt] at ($(W)+(0.05,-0.22)$) {$\dfrac{g\,t^a\,n_B^\nu}{n_B\cdot k-\ii0}$};
  \end{scope}
  % the exception: a single transverse gluon as the active parton (gamma* g -> q qbar)
  \begin{scope}[xshift=8.25cm,yshift=-5.85cm]
    \node[jet blob,ellipse,minimum width=1.75cm,minimum height=0.66cm,inner sep=0pt] (A) at (0,0) {};
    \node[font=\small] at (-0.42,0) {$A$};
    \node[hard blob,ellipse,minimum width=1.6cm,minimum height=0.66cm,inner sep=0pt] (H) at (0,2.25) {};
    \node[font=\small] at (H) {$H$};
    \draw[photon] (-1.25,3.05) -- (H.150);
    \draw[quark] (H.30) -- (1.25,3.0);
    \draw[quark] (1.35,1.75) -- (H.345);
    \draw[gluon] (0.05,0.33) -- (0.05,1.92);
    \path[mom'={$k$},mom sep=0.32cm] (0.05,0.33) -- (0.05,1.92);
    \node[fb idx,anchor=east] at (-0.05,1.15) {$j$};
    \node[fb idx,anchor=west] at (1.25,3.0) {$q$};
    \node[fb idx,anchor=west] at (1.35,1.75) {$\bar q$};
    \node[fb lab,align=center,anchor=north] (EXC) at (0,-0.45) {例外：活动部分子为单个横向胶子\\$k_\mu H^\mu=0$，G 项领头：\\$k^+A^j-k^jA^+\ \to\ G^{+j}$};
    \begin{scope}[on background layer]\node[fb box,fit={(EXC) (A) (H) (-1.3,3.1) (1.55,3.05)}] {};\end{scope}
  \end{scope}
\end{tikzpicture}
\caption{共线胶子的 Grammer--Yennie 分解（Feynman 规范；$k$ 由 $A$ 共线子图流入硬部分，$\mu$ 为硬端、$\nu$ 为共线端的指标）。
$g_{\mu\nu}=K_{\mu\nu}+G_{\mu\nu}$：K 部分（蓝色，圆点标出带标量极化 $k_\mu$ 的硬端）的共线端只取 $A^+=n_B\cdot A$ 分量；对 $H$ 内一切附着求和后，由 Ward 恒等式，K 胶子改为附着在沿 $n_B$ 的 Wilson 线上，乘以不含该胶子的硬部分 $H_0$（左下；$-\ii0$ 对应末态附着，即 DIS 中未来指向的线，初态附着取 $+\ii0$）；
G 部分满足 $G_{\mu\nu}k^\nu=0$，每多一条 G 胶子至少压低 $\lambda$。例外（右下）：作为活动部分子的单个横向胶子（此处为 $\gamma^*g\to q\bar q$），由 Ward 恒等式 $k_\mu H^\mu=0$，其 G 项是领头的，与 $k^+A^j-k^jA^+$ 组成场强 $G^{+j}$。}
\label{fig:method-kg}
\end{figure}
%%% END fig:method-kg

%%% BEGIN fig:method-neighbourhoods
\begin{figure}[htbp]
\centering
\begin{tikzpicture}[fd,
  fb lab/.style={font=\footnotesize},
  fb dim/.style={draw=black!75,line width=0.5pt,{Stealth[length=3.6pt,width=3pt]}-{Stealth[length=3.6pt,width=3pt]}},
  fb lead/.style={draw=black!60,line width=0.4pt}]
  \def\fbw{0.42}   % half-width lambda^{beta_C} of the collinear tubes
  \def\fbr{1.15}   % radius lambda^{beta_S} of the soft ball (beta_S < beta_C: fatter)
  % ---------------- the domain Omega and the dominance sets ----------------
  \begin{scope}
    \clip[rounded corners=9pt] (-3.75,-2.45) rectangle (3.75,2.45);
    \fill[gray!12] (-3.8,-2.5) rectangle (3.8,2.5);
    \fill[blue!15] (-3.8,-\fbw) rectangle (3.8,\fbw);
    \fill[orange!20] (-\fbw,-2.5) rectangle (\fbw,2.5);
    \draw[blue!55!black,line width=0.5pt,dash pattern=on 2pt off 1.5pt] (-3.8,-\fbw) -- (3.8,-\fbw) (-3.8,\fbw) -- (3.8,\fbw);
    \draw[orange!75!black,line width=0.5pt,dash pattern=on 2pt off 1.5pt] (-\fbw,-2.5) -- (-\fbw,2.5) (\fbw,-2.5) -- (\fbw,2.5);
    \fill[green!20] (0,0) circle[radius=\fbr];
    \draw[green!45!black,line width=0.6pt] (0,0) circle[radius=\fbr];
    \draw[blue!60!black,line width=1.1pt] (-3.8,0) -- (3.8,0);
    \draw[orange!80!black,line width=1.1pt] (0,-2.5) -- (0,2.5);
  \end{scope}
  \draw[gray!55!black,line width=0.7pt,rounded corners=9pt] (-3.75,-2.45) rectangle (3.75,2.45);
  \fill (0,0) circle[radius=1.8pt];
  % axes of the section k_T=0
  \draw[lcaxis,blue!60!black] (3.75,0) -- (4.1,0) node[right=0pt,font=\footnotesize,text=black] {$k^+$};
  \draw[lcaxis,orange!80!black] (0,2.45) -- (0,2.8) node[above=0pt,font=\footnotesize,text=black] {$k^-$};
  % pinch sets
  \node[fb lab,text=blue!50!black,anchor=south west] at (-3.62,\fbw+0.02) {$\pi_{C_1}$（$k\parallel p_1$）};
  \draw[fb lead,blue!50!black] (-3.15,\fbw+0.07) -- (-3.0,0.04);
  \node[fb lab,text=orange!60!black,anchor=north west] at (\fbw+0.06,2.36) {$\pi_{C_2}$（$k\parallel p_2$）};
  \draw[fb lead,orange!60!black] (\fbw+0.1,2.12) -- (0.04,1.95);
  \node[fb lab,anchor=south west,inner sep=1.5pt] at (0.05,0.03) {$\pi_S$};
  % dominance sets
  \node[font=\small,text=green!30!black] at (-0.66,-0.7) {$U_S$};
  \node[font=\small,text=blue!45!black] at (2.35,-0.21) {$U_{C_1}$};
  \node[font=\small,text=blue!45!black] at (-2.35,-0.21) {$U_{C_1}$};
  \node[font=\small,text=orange!55!black,fill=orange!20,inner sep=1pt] at (0,-1.78) {$U_{C_2}$};
  \node[font=\small,text=orange!55!black,fill=orange!20,inner sep=1pt] at (0,1.55) {$U_{C_2}$};
  \node[font=\small,text=gray!35!black] at (-2.35,1.55) {$U_H$};
  \node[font=\small,text=gray!35!black] at (2.35,-1.55) {$U_H$};
  \node[fb lab,anchor=south east,text=gray!30!black] at (3.68,-2.4) {$\Omega$};
  % dimensions
  \draw[fb dim] (0,0) -- (-45:\fbr);
  \draw[fb lead] (-45:0.62) -- (1.45,-1.12);
  \node[fb lab,anchor=west,inner sep=1pt] at (1.45,-1.12) {$\lambda^{\beta_S}$};
  \draw[fb dim] (3.05,0) -- (3.05,\fbw);
  \draw[fb lead] (3.05,0.3) -- (3.3,0.75);
  \node[fb lab,anchor=south,inner sep=1pt] at (3.3,0.75) {$\lambda^{\beta_C}$};
  % ---------------- definitions ----------------
  \node[anchor=north west,align=left,font=\footnotesize,text width=6.3cm] at (4.75,2.42) {%
    $N_R=\{\sigma_{R(j)}(l_j)\le\lambda^{\beta_R}\ \ \forall j:\ R(j)\neq H\}$\\[2pt]
    $\sigma_A=|l^-|+|\mathbf l_T|$，\ $\sigma_B=|l^+|+|\mathbf l_T|$，\ $\sigma_S=\|l\|$\\[6pt]
    $U_R=N_R\setminus\bigcup_{R'<R}N_{R'}$：\\[2pt]
    \quad$U_S=N_S$，\ \ $U_{C_i}=N_{C_i}\setminus N_S$，\\[2pt]
    \quad$U_H=\Omega\setminus(N_S\cup N_{C_1}\cup N_{C_2})$\\[6pt]
    $\beta_R=3^{h(R)-h_*}$：较小区域的邻域较“胖”\\[2pt]
    （本例 $\beta_S=\frac19<\beta_C=\frac13$，$\lambda^{\beta_S}>\lambda^{\beta_C}$）\\[6pt]
    $\bigcup_RU_R=\Omega$（\cref{lem:method-covering}）};
\end{tikzpicture}
\caption{依赖于 $\lambda$ 的邻域 $N_R=\{\sigma_{R(j)}(l_j)\le\lambda^{\beta_R}\}$ 与支配集 $U_R=N_R\setminus\bigcup_{R'<R}N_{R'}$（\cref{def:method-neighbourhoods}；示意：单圈离壳三角图的胶子动量 $k$ 在 $\mathbf k_T=0$ 处的 $(k^+,k^-)$ 截面，以 $Q$ 为单位）。
共线钳夹集 $\pi_{C_1}$、$\pi_{C_2}$ 相交于软钳夹点 $\pi_S$（$k=0$）；$N_{C_i}$ 是半宽 $\lambda^{\beta_C}$ 的管道，$N_S$ 是半径 $\lambda^{\beta_S}$ 的球，
由于 $\beta_S<\beta_C$ 而比管道“胖”。绿、蓝、橙、灰色区域分别为 $U_S$、$U_{C_1}$、$U_{C_2}$、$U_H$，它们覆盖整个积分域（\cref{lem:method-covering}）。}
\label{fig:method-neighbourhoods}
\end{figure}
%%% END fig:method-neighbourhoods

%%% BEGIN fig:method-nesting
\begin{figure}[htbp]
\centering
\begin{tikzpicture}[fd,
  fb lab/.style={font=\footnotesize},
  fb ellA/.style={rotate around={22:(-0.95,0)}},
  fb ellB/.style={rotate around={-22:(0.95,0)}},
  fb step/.style={draw=gray!40!black,line width=0.75pt,-{Stealth[length=5pt,width=4.2pt]}},
  fb stepl/.style={font=\footnotesize,inner sep=1.5pt,fill=white,rounded corners=1.5pt}]
  % ---------------- Venn-like diagram of the pinch sets (schematic) ----------------
  \draw[fill=gray!10,draw=gray!55!black,line width=0.7pt,rounded corners=10pt] (-3.6,-2.35) rectangle (3.6,2.35);
  \fill[blue!12,fb ellA] (-0.95,0) ellipse[x radius=2.05,y radius=1.08];
  \fill[orange!17,fb ellB] (0.95,0) ellipse[x radius=2.05,y radius=1.08];
  % pi_S = pi_A cap pi_B: the lens common to both ellipses
  \begin{scope}
    \clip[fb ellA] (-0.95,0) ellipse[x radius=2.05,y radius=1.08];
    \fill[green!22,fb ellB] (0.95,0) ellipse[x radius=2.05,y radius=1.08];
  \end{scope}
  \draw[blue!55!black,line width=0.7pt,fb ellA] (-0.95,0) ellipse[x radius=2.05,y radius=1.08];
  \draw[orange!75!black,line width=0.7pt,fb ellB] (0.95,0) ellipse[x radius=2.05,y radius=1.08];
  \node[font=\small,text=green!30!black] at (0,0.22) {$\pi_S$};
  \node[font=\small,text=blue!50!black] at (-2.05,-0.25) {$\pi_A$};
  \node[font=\small,text=orange!60!black] at (2.05,-0.25) {$\pi_B$};
  \node[font=\small,anchor=north west] at (-3.5,2.25) {$\pi_H=\Omega$};
  \node[fb lab,text=gray!30!black,anchor=south] at (0,-2.3) {$\pi_S=\pi_A\cap\pi_B$（\cref{lem:method-meets}）};
  % order of the approximations: smallest region first
  \draw[fb step] (0.1,1.95) node[fb stepl,anchor=south] {\textcircled{\scriptsize 1}\,$T_S$} to[out=-90,in=90] (0.04,0.62);
  \draw[fb step] (-2.9,-1.65) node[fb stepl,anchor=north] {\textcircled{\scriptsize 2}\,$T_A$} to[out=60,in=-150] (-1.85,-0.75);
  \draw[fb step] (2.9,-1.65) node[fb stepl,anchor=north] {\textcircled{\scriptsize 2}\,$T_B$} to[out=120,in=-30] (1.85,-0.75);
  \draw[fb step] (2.75,1.85) node[fb stepl,anchor=south west,xshift=-0.15cm] {\textcircled{\scriptsize 3}\,$T_H$} to[out=-120,in=60] (2.4,1.2);
  % ---------------- subtraction terms ----------------
  \begin{scope}[xshift=4.15cm]
    \node[anchor=north west,align=left,font=\small] (F) at (0,2.2) {%
      $C_S\Gamma=T_S\,\Gamma$\\[4pt]
      $C_A\Gamma=T_A(1-T_S)\,\Gamma$，\quad$C_B\Gamma=T_B(1-T_S)\,\Gamma$\\[4pt]
      $C_H\Gamma=T_H\big(\Gamma-C_S\Gamma-C_A\Gamma-C_B\Gamma\big)$\\[4pt]
      $\Gamma=\sum_{R}C_R\Gamma+r$};
    \node[anchor=north west,align=left,font=\footnotesize,text width=6.6cm,draw=red!50!black,
          line width=0.5pt,rounded corners=2pt,fill=red!3,inner sep=4pt] at ($(F.south west)+(0.05,-0.25)$)
      {只有分组表达式被压低（\cref{ex:method-naive-induction}）：单个链项一般不可积（例如离壳三角图中的 $T_{C_1}\Gamma$ 对任何 $\varepsilon$ 都不绝对可积），被压低并满足 \mbox{(CS1)--(CS3)} 的是分组表达式 $C_R\Gamma$ 与余项 $r$。};
  \end{scope}
\end{tikzpicture}
\caption{区域的偏序与减除（示意；钳夹集画成区域）。单圈顶点的钳夹集彼此嵌套：软钳夹集 $\pi_S$（绿色）同时含于两个共线钳夹集 $\pi_A$（蓝）与 $\pi_B$（橙），并且恰为二者之交 $\pi_S=\pi_A\cap\pi_B$，即逐线交 $A\wedge B=S$；所有钳夹集都含于硬区 $\pi_H=\Omega$。
减除自最小区域（软）开始，依次为 \textcircled{\scriptsize 1}~$T_S$、\textcircled{\scriptsize 2}~$T_A,T_B$、\textcircled{\scriptsize 3}~$T_H$：较大区域的近似子作用于减去较小区域贡献之后的被积函数，$C_R\Gamma\equiv T_R\big(\Gamma-\sum_{R'<R}C_{R'}\Gamma\big)$（\cref{def:method-subtraction}）。
单个链项一般不可积，只有分组表达式满足 \mbox{(CS1)--(CS3)}。}
\label{fig:method-nesting}
\end{figure}
%%% END fig:method-nesting

%%% BEGIN fig:method-sudakov-plane
\begin{figure}[htbp]
\centering
\begin{tikzpicture}[fd,
  fb lab/.style={font=\footnotesize,inner sep=1.5pt},
  fb tick/.style={font=\scriptsize},
  fb grid/.style={draw=gray!45,line width=0.4pt,dash pattern=on 1pt off 2pt},
  fb hyp/.style={draw=black!60,line width=0.6pt},
  fb patch/.style={circle,draw=black!70,line width=0.6pt,minimum size=0.62cm,inner sep=0pt,font=\small},
  fb pole/.pic={\draw[line width=1pt,line cap=round] (-0.085,-0.085)--(0.085,0.085) (-0.085,0.085)--(0.085,-0.085);},
  fb contour/.style={draw=blue!60!black,line width=1.0pt},
  fb carrow/.style={postaction={decorate,decoration={markings,
      mark=at position #1 with {\arrow{Stealth[length=5.5pt,width=4.5pt]}}}}}]
  % ======================= (a) the (ln|k^+|, ln|k^-|) plane =======================
  \begin{scope}[x=1.15cm,y=1.15cm]
    % soft--collinear overlap strips (hatched)
    \fill[pattern=north east lines,pattern color=green!45!black] (-2,-2.15) rectangle (0,-1.85);
    \fill[pattern=north east lines,pattern color=green!45!black] (-2.15,-2) rectangle (-1.85,0);
    % axes, ticks, grid
    \draw[lcaxis] (-2.85,-2.85) -- (1.1,-2.85) node[below=1pt,font=\small] {$\ln|k^+|$};
    \draw[lcaxis] (-2.85,-2.85) -- (-2.85,1.1) node[left=1pt,font=\small] {$\ln|k^-|$};
    \foreach \v/\t in {-2/{\lambda^2Q},-1/{\lambda Q},0/{Q}} {
      \draw[line width=0.5pt] (\v,-2.85) -- ++(0,0.07);
      \node[fb tick,below] at (\v,-2.85) {$\t$};
      \draw[line width=0.5pt] (-2.85,\v) -- ++(0.07,0);
      \node[fb tick,left] at (-2.85,\v) {$\t$};
      \draw[fb grid] (\v,-2.78) -- (\v,0.95);
      \draw[fb grid] (-2.78,\v) -- (0.95,\v);
    }
    % lines of constant k^+ k^-
    \draw[fb hyp] (0.75,-0.75) -- (-0.98,0.98) node[left=1pt,fb lab,text=black!70] {$k^+k^-\sim Q^2$};
    \draw[fb hyp,line width=1.0pt] (-2.62,0.62) -- (0.62,-2.62) node[right=1pt,fb lab,text=black!85] {$k^+k^-\sim m^2$};
    % the family k=(z, lambda^2/z, lambda) between C1 and C2
    \draw[red!60!black,line width=1.6pt,dash pattern=on 4pt off 2pt] (-1.72,-0.28) -- (-0.28,-1.72);
    % regions of the off-shell triangle
    \node[fb patch,fill=gray!35] (H) at (0,0) {$H$};
    \node[fb patch,fill=blue!20] (C1) at (0,-2) {$C_1$};
    \node[fb patch,fill=orange!25] (C2) at (-2,0) {$C_2$};
    \node[fb patch,fill=green!25] (S) at (-2,-2) {$S$};
    \draw[black!70,line width=0.6pt,fill=white] (-1,-1) circle[radius=2.3pt];
    \draw[black!60,line width=0.4pt] (-1.06,-1.06) -- (-1.3,-1.4);
    \node[font=\scriptsize,inner sep=1pt,anchor=north] at (-1.36,-1.4) {$(\lambda,\lambda,\lambda)$};
    \node[fb lab,text=red!55!black,anchor=north west,align=left] at (0.02,-0.97) {$k=(z,\lambda^2/z,\lambda)$\\$\lambda^2\ll z\ll1$：见 (b)};
    \node[fb lab,text=green!30!black] at (-1,-2.42) {软—共线重叠};
    \node[font=\small,anchor=south west] at (-3.75,1.15) {(a)};
  \end{scope}
  % ======================= (b) complex k^- plane =======================
  \begin{scope}[xshift=5.55cm,yshift=-1.4cm]
    \fill[red!12] (0,0) circle[radius=0.42];
    \draw[lcaxis] (-2.3,0) -- (3.2,0) node[right=1pt,font=\small] {$\Real k^-$};
    \draw[lcaxis] (0,-1.25) -- (0,1.85) node[left=1pt,font=\small] {$\Imag k^-$};
    % contour: lower half plane on the left, crossing near k_T^2/(4k^+), upper half plane on the right
    \draw[fb contour,fb carrow=0.18,fb carrow=0.86] (-2.2,-1.0) .. controls (-0.6,-1.0) and (0.55,-0.55) .. (1.1,0)
         .. controls (1.6,0.45) and (2.25,0.95) .. (3.05,0.95);
    \pic at (-0.22,0.2) {fb pole};
    \pic at (2.2,-0.2) {fb pole};
    \fill[blue!60!black] (1.1,0) circle[radius=1.5pt];
    \node[fb lab,anchor=south] at (-0.22,0.3) {$D_1$};
    \node[fb lab,anchor=north] at (2.2,-0.3) {$D_g$};
    \draw[gray!50!black,line width=0.4pt] (1.12,0.07) -- (1.28,1.18) node[above,fb lab] {$\Real k^-\approx\dfrac{k_T^2}{4k^+}$};
    \draw[red!50!black,line width=0.4pt] (-0.32,0.28) -- (-0.62,0.56);
    \node[fb lab,text=red!55!black,anchor=south east,align=right] at (-0.58,0.5) {$k^-\approx0$：\\软近似失效};
    \node[fb lab,anchor=north west,align=left] at (-2.2,-1.42) {$D_1$：$k^-\approx-\dfrac{m^2+k_T^2}{2P}+\ii0$\\[2pt]$D_g$：$k^-=\dfrac{k_T^2-\ii0}{2k^+}$};
    \node[font=\small,anchor=south west] at (-2.3,1.55) {(b)};
  \end{scope}
\end{tikzpicture}
\caption{单圈离壳三角图（$p_i^2=-m^2$，$\lambda=m/Q$）的区域与 Collins 减除中的软—共线重叠（\cref{ex:method-sudakov-subtraction}）。
(a)~$(\ln|k^+|,\ln|k^-|)$ 平面（以 $\ln(1/\lambda)$ 为单位，$k_T$ 取 $k_T^2\sim|k^+k^-|$ 的典型值）：区域 $H$、$C_1$、$C_2$、$S$ 位于正方形的四角，其中软区 $S\sim Q(\lambda^2,\lambda^2,\lambda^2)$ 是\cref{def:ir-scalings}的超软标度（$k\sim m^2/Q$）；
粗斜线 $k^+k^-\sim m^2$ 是 $k_T\sim\lambda Q$ 的在壳线 $k_T^2=2k^+k^-$，经过 $C_1$、$C_2$ 与“软中心”$(\lambda,\lambda,\lambda)$；绿色阴影带为软—共线重叠，减除项在重叠区相消的方式由分组表达式决定。
(b)~对 $k=(z,\lambda^2/z,\lambda)$（$\lambda^2\ll z\ll1$，红色虚线），胶子线 $D_g$ 与夸克线 $D_1$ 的 $k^-$ 极点分居实轴两侧、相距 $\sim\lambda^2Q/z$，围道只能在两极点之间穿过实轴；
软近似只在 $\Real k^-\sim k_T^2/(4k^+)$ 附近穿过实轴的围道上成立，在 $k^-\approx0$ 附近（红色区域）失效（\cref{lem:method-i0-compat}）。}
\label{fig:method-sudakov-plane}
\end{figure}
%%% END fig:method-sudakov-plane

%%% BEGIN fig:ward-identity
\begin{figure}[htbp]
\centering
\begin{tikzpicture}[fd,
  t2 lab/.style={font=\small},
  t2 eq/.style={font=\large},
  t2 hb/.style={hard blob,ellipse,minimum width=1.3cm,minimum height=0.95cm,inner sep=0pt}]
  % ======================= row (a): one K gluon =======================
  \node[t2 eq] at (-2.35,-0.1) {$\displaystyle\sum_{\text{所有附着}}$};
  \node[t2 hb] (H) at (0,0) {$H$};
  \draw[photon] (-1.25,1.05) -- (H.140);
  \draw[quark] (H.25) -- (1.55,0.85);
  \draw[quark] (-1.5,-1.05) -- (H.214);
  \path[mom={$p$},mom pos=0.4] (-1.5,-1.05) -- (H.214);
  \draw[longit] (-0.55,-1.6) -- (H.264);
  \longitmark{(-0.55,-1.6)}{(H.264)}
  \path[mom'={$k$},mom sep=0.3cm] (-0.55,-1.6) -- (H.264);
  \node[t2 lab,anchor=north] at (-0.55,-1.6) {$a$};
  \node[t2 eq] at (2.35,0) {$=$};
  \begin{scope}[xshift=4.6cm]
    \node[t2 hb] (H0) at (0,0) {$H_0$};
    \draw[photon] (-1.25,1.05) -- (H0.140);
    \draw[quark] (H0.25) -- (1.55,0.85);
    \draw[quark] (-1.5,-1.05) -- (H0.214);
    \path[mom={$p+k$},mom pos=0.4] (-1.5,-1.05) -- (H0.214);
    \draw[wilson] (H0.214) -- ++(-33:2.25) coordinate (W) node[right,t2 lab] {沿 $n_B$};
    \coordinate (Y) at ($(H0.214)!0.52!(W)$);
    \draw[longit] (-0.55,-1.6) -- (Y);
    \longitmark{(-0.55,-1.6)}{(Y)}
    \path[mom'={$k$},mom sep=0.3cm] (-0.55,-1.6) -- (Y);
    \node[t2 lab,anchor=north] at (-0.55,-1.6) {$a$};
  \end{scope}
  \node[font=\small,anchor=west] at (-2.95,-2.4) {(a)\ 阿贝尔情形（任意多个 K 光子），或单个 K 胶子};
  % ======================= row (b): two K gluons, non-abelian =======================
  \begin{scope}[yshift=-4.2cm]
    \node[t2 eq] at (-2.35,-0.1) {$\displaystyle\sum_{\text{所有附着}}$};
    \node[t2 hb] (H) at (0,0) {$H$};
    \draw[photon] (-1.25,1.05) -- (H.140);
    \draw[quark] (H.25) -- (1.55,0.85);
    \draw[quark] (-1.5,-1.05) -- (H.214);
    \draw[longit] (-0.75,-1.62) -- (H.250);  \longitmark{(-0.75,-1.62)}{(H.250)}
    \draw[longit] (0.05,-1.68) -- (H.290);   \longitmark{(0.05,-1.68)}{(H.290)}
    \node[t2 lab,anchor=north] at (-0.75,-1.62) {$k_1$};
    \node[t2 lab,anchor=north] at (0.05,-1.68) {$k_2$};
    \node[t2 eq] at (2.35,0) {$=$};
    \begin{scope}[xshift=4.6cm]
      \node[t2 hb] (H0) at (0,0) {$H_0$};
      \draw[photon] (-1.25,1.05) -- (H0.140);
      \draw[quark] (H0.25) -- (1.55,0.85);
      \draw[quark] (-1.5,-1.05) -- (H0.214);
      \draw[wilson=0.84] (H0.214) -- ++(-33:2.25) coordinate (W) node[right,t2 lab] {沿 $n_B$};
      \draw[longit] (-0.75,-1.62) -- ($(H0.214)!0.36!(W)$);  \longitmark{(-0.75,-1.62)}{($(H0.214)!0.36!(W)$)}
      \draw[longit] (0.05,-1.68) -- ($(H0.214)!0.66!(W)$);   \longitmark{(0.05,-1.68)}{($(H0.214)!0.66!(W)$)}
      \node[t2 lab,anchor=north] at (-0.75,-1.62) {$k_1$};
      \node[t2 lab,anchor=north] at (0.05,-1.68) {$k_2$};
      \node[font=\small,anchor=west] at (1.7,-0.08) {$+\;(k_1\leftrightarrow k_2)$};
    \end{scope}
    \node[t2 eq] at (8.75,0) {$+$};
    \begin{scope}[xshift=10.55cm]
      \node[t2 hb] (Hg) at (0,0) {$H'$};
      \draw[photon] (-1.25,1.05) -- (Hg.140);
      \draw[quark] (Hg.25) -- (1.55,0.85);
      \draw[quark] (-1.5,-1.05) -- (Hg.214);
      \draw[ghost] (-0.85,-1.65) -- (Hg.252);
      \draw[ghost] (Hg.292) -- (0.15,-1.7);
      \node[font=\footnotesize,anchor=north,align=center] at (-0.35,-1.75) {鬼场线\\（鬼场项，示意）};
    \end{scope}
  \end{scope}
  \node[font=\small,anchor=west] at (-2.95,-6.85) {(b)\ 非阿贝尔情形，两个及以上 K 胶子};
\end{tikzpicture}
\caption{图示 Ward 恒等式。蓝色螺旋线是标量（纵向）极化的共线胶子（K 胶子，极化 $\propto k^\mu$，以圆点标记），来自 $A$ 共线子图（未画出）；
$H$ 的其余外线（入射夸克、出射夸克、光子）都取为在壳的物理粒子（共线近似之后）。
(a)~对 K 胶子在 $H$ 内一切可能附着位置（包括非阿贝尔情形中 $H$ 内的胶子线与鬼场线）求和，等于把它附着到一条沿 $n_B$ 方向的 Wilson 线（双线）上、
再乘以不含该胶子的硬部分 $H_0$（其入射动量为 $p+k$）；Wilson 线是未来指向还是过去指向，取决于 K 胶子附着在 $H$ 的末态线还是初态线上。
阿贝尔理论中这一等式对任意多个 K 光子成立；非阿贝尔理论中对单个 K 胶子成立（BRST 对称性）。
(b)~非阿贝尔理论中有两个及以上 K 胶子时，求和给出按路径次序附着在 Wilson 线上的各项（色矩阵 $t^{a_1}$、$t^{a_2}$ 的次序由附着次序决定，两种次序都出现），
一般还伴随含鬼场线（带箭头的点线）的附加项；此处只示意其存在，其精确形式与处理见附录 F。}
\label{fig:ward-identity}
\end{figure}
%%% END fig:ward-identity
